\begin{theorem}[window-\texorpdfstring{$6$}{6} 的端点原子通道不能经有限 \texorpdfstring{$2$}{2}-主异常账本加性实现]\label{thm:conclusion-window6-endpoint-atom-not-additive-through-2primary-ledger}
\leanverified{paper\_conclusion\_window6\_endpoint\_atom\_not\_additive\_through\_2primary\_ledger}
沿用定理 \ref{thm:conclusion-window6-boundary-parity-directsummand-rational-blindness} 与推论 \ref{cor:conclusion-window6-boundary-parity-misses-eighteen-anomaly-directions} 的记号。则任意群同态
\[
\ell:\bigl(G_6^{\mathrm{bin}}\bigr)^{\mathrm{ab}}\longrightarrow (\mathbb R,+)
\]
都必为零。因而任何只经由有限 \(2\)-主异常账本
\[
\bigl(G_6^{\mathrm{bin}}\bigr)^{\mathrm{ab}}
\cong
(\ZZ/2\ZZ)^3\oplus(\ZZ/2\ZZ)^{18}
\]
因子化的实值加性可观测量都只能是常值。

特别地，端点原子质量通道不能由该异常账本的加性像忠实实现；若要同时读取 boundary parity、隐藏异常类与端点原子载荷，则必须在有限 \(2\)-主账本之外额外并列一条非挠实标量轴。
\end{theorem}
\begin{proof}
由定理 \ref{thm:conclusion-window6-boundary-parity-directsummand-rational-blindness}，
\[
\bigl(G_6^{\mathrm{bin}}\bigr)^{\mathrm{ab}}
\cong
C_{\partial}\oplus C_{\mathrm{bulk}}
\cong
(\ZZ/2\ZZ)^3\oplus(\ZZ/2\ZZ)^{18}
\]
是有限 \(2\)-群。另一方面，\((\mathbb R,+)\) 无非零挠元，故任意群同态
\(
\ell:(G_6^{\mathrm{bin}})^{\mathrm{ab}}\to (\mathbb R,+)
\)
都必为零。

于是，一切只经由该 abelian anomaly ledger 因子化的实值加性可观测量都只能是常值。可另一方面，正文关于端点热核探针与端点原子读出的结果已表明：端点原子质量是一条独立的边界层通道，并且不能由其余可见/离散通道消去。故若要忠实读取端点原子载荷，它就不可能仅靠有限 \(2\)-主异常账本的加性像来承载，而必须额外附加一条非挠实标量轴。证毕。
\end{proof}

\begin{conjecture}[隐藏异常商到端点奇异层的 transgression]\label{conj:conclusion-window6-hidden-anomaly-to-endpoint-transgression}
设
\[
\bigl(G_6^{\mathrm{bin}}\bigr)^{\mathrm{ab}}/C_{\partial}
\cong
(\ZZ/2\ZZ)^{18}
\]
为 window-\(6\) 的隐藏异常商，并记端点原子通道中的奇异层为 \(\mathcal E_{\partial}^{\mathrm{sing}}\)。则应存在一个自然的 transgression 映射
\[
\mathrm{tr}_{\partial}:
(\ZZ/2\ZZ)^{18}\longrightarrow \mathcal E_{\partial}^{\mathrm{sing}},
\]
使得：
\begin{enumerate}
  \item \(\ker(\mathrm{tr}_{\partial})\) 恰刻画那些在隐藏异常账本中非平凡但在端点原子读出上仍不可见的类；
  \item \(\operatorname{im}(\mathrm{tr}_{\partial})\) 正是端点原子通道中无法由 visible modulus 与 interior-zero 两轴解释的最小外置部分；
  \item \(\mathrm{tr}_{\partial}=0\) 当且仅当边界审计可由模长轴、盘内零点轴与有限 \(2\)-主异常账本闭合，而不再需要额外实轴。
\end{enumerate}
\end{conjecture}
